Long-horizon agents must acquire resources, finance plans and negotiate while earlier actions reshape later opportunities. We study economic capability and the social conduct of bargaining in MonopolyBench, a deterministic four-player environment with recorded legal menus, response attempts, delivered memory and economic events. Across twenty completed games, we analyze 6,989 decisions and 8,679 attempts using event reconstruction, overlapping source review and exact-prefix execution of local alternatives. Forty-three of sixty-three trade-completed groups receive construction in the first management window. Among forty-two groups buildable on entry, eight remain unbuilt in that window: three through engine defaults, three through received actions conflicting with accompanying content, and two through stated holds. Three of sixty eliminations admit a legal current-debt payment sequence; another has sufficient nominal value but is constrained by hotel liquidation. Social chronologies show a misleading omission, recipient correction alongside acceptance, a transient false echo, and withdrawal of a standing offer. Some received proposals contradict their accompanying text, so settlement alone identifies neither persuasion nor coherent proposer intent. All twenty recorded action streams replay, and all 6,989 memory windows reconstruct. Public deception is permitted, and one configuration operates almost entirely under a provider-refusal recovery condition without authored private notes. The cohort therefore establishes neither an instruction-misalignment rate nor a capability--alignment trend. It contributes an empirical account, an exact-prefix audit procedure applied to agents' claims and research labels, and requirements for future controlled, norm-explicit comparisons.
